\title{\textbf{Applied Vacuum Engineering} \\ \large \textit{Volume III: The Macroscopic Continuum}}
\author{Grant Lindblom}
\date{}

\maketitle

\vfill
\noindent\textbf{Applied Vacuum Engineering: Volume III} \\
This document unites microscopic topological soliton structure with macroscopic thermodynamics, gravity, and cosmology.

\begin{abstract}
    \textbf{Volume III: The Macroscopic Continuum} applies the substrate model across scales, treating gravitation as the macroscopic scalar impedance gradient ($Z$) of the LC network rather than as an independent geometric manifold. The stress-energy content is the electromagnetic energy density of the local LC vacuum:
    \begin{equation}
        T_{\mu\nu} \equiv U_{\mu\nu} = \tfrac{1}{2}\epsilon_0 |\mathbf{E}|^2 + \tfrac{1}{2}\mu_0 |\mathbf{H}|^2
    \end{equation}
    the volume develops gravity, thermodynamic entropy, superconductivity, and the universe's expansion limits (dark energy) from the substrate's tension and the latent heat of ongoing lattice crystallization. The gravitational \emph{form} is substrate-derived; $G$'s value is treated as a calibration input (see the common foreword's honest-scope notes).
\end{abstract}
